\documentclass[11pt]{article}
\usepackage[a4paper,margin=24mm]{geometry}
\usepackage{amsmath,amssymb,amsthm,mathtools,bm}
\usepackage{booktabs,array,microtype,xurl,hyperref}
\hypersetup{colorlinks=true,linkcolor=blue,citecolor=blue,urlcolor=blue}

\newtheorem{theorem}{Theorem}[section]
\newtheorem{proposition}[theorem]{Proposition}
\newtheorem{lemma}[theorem]{Lemma}
\newtheorem{corollary}[theorem]{Corollary}
\newtheorem{definition}[theorem]{Definition}
\newcommand{\dd}{\,\mathrm d}
\newcommand{\RR}{\mathbb R}
\newcommand{\ZZ}{\mathbb Z}
\newcommand{\cuph}{\smile_h}
\newcommand{\Wh}{\mathcal W_h}
\newcommand{\Ih}{I_h}
\newcommand{\tr}{\operatorname{tr}}
\newcommand{\cof}{\operatorname{cof}}

\title{AP-E11 Compatible Tetrahedral Cochain Action:\
Exact Leibniz Algebra, Positive Whitney Hodge Star, Shared Degree Cochain,\
and a Mesh-Independent Complete-Minor Cell Formula}
\author{Route-F derivation ledger}
\date{19 July 2026}

\begin{document}
\maketitle

\begin{abstract}
AP-E10 proved that the one-corner forward Skyrme stencil is not a compatible
discrete differential product.  AP-E11 replaces it by one ordered
tetrahedral cochain complex.  The Alexander--Whitney front/back cup satisfies
the exact Leibniz identity; its target-antisymmetric products are precisely
the de Rham cochains of the affine two- and three-minors.  Consistent Whitney
mass matrices give positive, rigid-rotation-covariant Hodge stars.  A radial
dressing of the same third-cup family is the exact degree density of the
normalized piecewise-affine map.  Adding the positive three-cochain Hodge
norm closes the finite-grid collapse mode and gives the polyconvex density
\[
 Q_{R,K}(A)=\frac12|A|^2+\frac R2|\wedge^2A|^2
              +\frac K2|\wedge^3A|^2 .
\]
Every periodic P1 corrector preserves the means of all three minor orders;
Jensen's inequality therefore proves that the corrector is zero and the cell
formula is identical on the six- and checkerboard five-tetrahedron meshes.
The production card passes $11/11$ exact/mechanical checks, all thirteen
unanchored $B=1$ solves, and the declared translation/triangulation quotient
continuum gate.  What is closed unconditionally is Gamma convergence to the
fixed-degree lower-semicontinuous relaxation.  Identification of that
relaxation with the unrelaxed classical graph-norm functional remains a
separate density theorem; consequently the classical Hessian and determinant
variation are not yet authorized.
\end{abstract}

\tableofcontents

\section{Why the action must be replaced}

Let $K_h$ be an ordered, oriented, shape-regular tetrahedral complex in a
bounded Lipschitz domain $\Omega\subset\RR^3$.  A nodal field is a
zero-cochain $n_h=(n_h^0,\ldots,n_h^3)\in C^0(K_h;\RR^4)$ with
$|n_h(v)|=1$ at every vertex.  AP-E10 sampled all spatial derivatives from
one cube corner and multiplied them without a cochain shift.  The exact
oscillatory counterexample then retained a nonzero discrete two-minor while
converging uniformly to a constant.  No coefficient calibration can repair
that algebraic failure.

The replacement has three requirements.
First, $d^2=0$ and the product used to form a minor must satisfy a discrete
Leibniz rule.  Second, the positive quadratic form called a Hodge star must
reconstruct affine exterior forms independently of a cube body diagonal.
Third, the topological three-cochain must refer to the same normalized-affine
map as the finite-grid degree.  Whitney forms and finite element exterior
calculus provide the appropriate commuting-complex architecture
\cite{Whitney1957,Dodziuk1976,ArnoldFalkWinther2010}.  Recent discrete
exterior-calculus work likewise emphasizes that a proposed wedge must be
checked against Leibniz rather than inferred from visual locality
\cite{PtackovaVelho2024,Ptackova2025}; those surface results are motivation,
not a proof of the present three-dimensional construction.

\section{Exact shifted cup algebra}

For an increasing simplex $[v_0\ldots v_{p+q}]$, define the front/back
Alexander--Whitney product
\begin{equation}
 (\alpha\cuph\beta)[v_0\ldots v_{p+q}]
 =\alpha[v_0\ldots v_p],\beta[v_p\ldots v_{p+q}],
 \qquad \alpha\in C^p,\quad\beta\in C^q .
 \label{eq:aw}
\end{equation}

\begin{theorem}[Differential graded identities]
The simplicial coboundary and \eqref{eq:aw} obey
\begin{align}
 d^2&=0,\label{eq:d2}\\
 d(\alpha\cuph\beta)
 &=d\alpha\cuph\beta+(-1)^p\alpha\cuph d\beta,
 \label{eq:leibniz}\\
 (\alpha\cuph\beta)\cuph\gamma
 &=\alpha\cuph(\beta\cuph\gamma).
 \label{eq:assoc}
\end{align}
\end{theorem}

\begin{proof}
Expand $d$ as the alternating sum of deleted vertices.  Deleting one of the
front vertices produces the first term in \eqref{eq:leibniz}; deleting a
back vertex produces the second, including the sign from passing $p$
front degrees.  The two copies obtained by deleting the shared vertex cancel.
A second deletion cancels in pairs, proving \eqref{eq:d2}.  Both sides of
\eqref{eq:assoc} evaluate to the product on the three consecutive front,
middle, and back blocks.
\end{proof}

Put $a^A=dn_h^A$.  The correctly normalized antisymmetric two-cochain is
\begin{equation}
 X_h^{AB}=\frac1{2!}\left(a^A\cuph a^B-a^B\cuph a^A\right).
 \label{eq:two-cup}
\end{equation}
On an oriented face $[012]$,
\begin{equation}
 X_h^{AB}[012]
 =\frac12\det
 \begin{pmatrix}
 n_1^A-n_0^A & n_2^A-n_0^A\\
 n_1^B-n_0^B & n_2^B-n_0^B
 \end{pmatrix}
 =\int_{[012]}\dd(\Ih n^A)\wedge\dd(\Ih n^B).
 \label{eq:two-exact}
\end{equation}
The factor $1/2!$ is essential: the reference triangle has area $1/2$.
Similarly define the fully antisymmetrized family
\begin{equation}
 Z_h^{ABC}=\frac1{3!}\sum_{\pi\in S_3}\operatorname{sgn}(\pi)
 a^{\pi(A)}\cuph a^{\pi(B)}\cuph a^{\pi(C)}.
 \label{eq:three-cup}
\end{equation}
It is the de Rham three-cochain of the affine third minor.

There is also a distinguished target contraction.  If
$T=[0123]$, direct expansion gives
\begin{equation}
 C_h[n](T)=\epsilon_{ABCD}
 \left(n^A\cuph a^B\cuph a^C\cuph a^D\right)(T)
 =\det[n_0\ n_1\ n_2\ n_3].
 \label{eq:det-cup}
\end{equation}
Thus the energy cochains \eqref{eq:three-cup} and the topological contraction
\eqref{eq:det-cup} belong to one differential graded algebra; they are not
separate stencils.

\section{Positive rotation-calibrated Hodge stars}

Let $\lambda_i$ be the barycentric coordinates of a tetrahedron.  For an
oriented $k$-face $\sigma=[i_0\ldots i_k]$, the Whitney form is
\begin{equation}
 \Wh(\sigma)=k!\sum_{r=0}^k(-1)^r\lambda_{i_r}
 \dd\lambda_{i_0}\wedge\cdots\widehat{\dd\lambda_{i_r}}
 \cdots\wedge\dd\lambda_{i_k}.
 \label{eq:whitney}
\end{equation}
Define the local consistent Hodge matrix by
\begin{equation}
 (M_k^T)_{\sigma\tau}
 =\int_T\Wh(\sigma)\wedge\star\Wh(\tau),
 \qquad k=1,2,3.
 \label{eq:hodge}
\end{equation}

\begin{proposition}[Positivity and rotation calibration]
For every nondegenerate tetrahedron, $M_k^T$ is positive definite.  Under a
rigid rotation $Q\in SO(3)$, the naturally identified matrix satisfies
$M_k^{QT}=M_k^T$.  If $f,g$ are affine, then
\begin{align}
 (df)^TM_1^T(df)&=\int_T|\dd f|^2,\label{eq:hodge1}\\
 (X_h^{fg})^TM_2^T(X_h^{fg})&=\int_T|\dd f\wedge\dd g|^2,
 \label{eq:hodge2}\\
 (Z_h^{fgh})^TM_3^T(Z_h^{fgh})&=\int_T|\dd f\wedge\dd g\wedge\dd h|^2.
 \label{eq:hodge3}
\end{align}
\end{proposition}

\begin{proof}
Equation \eqref{eq:hodge} is a Gram matrix.  Linear independence of the local
Whitney basis makes its nullspace trivial.  A rigid rotation preserves volume
and the Euclidean inner product on exterior forms, giving covariance.  The
commuting identity $\Wh d=\dd\Wh$ reconstructs every constant differential
of an affine function.  Equations \eqref{eq:two-exact} and
\eqref{eq:three-cup} then give \eqref{eq:hodge1}--\eqref{eq:hodge3}.
\end{proof}

This consistent star is deliberately not replaced by a diagonal
circumcentric formula; diagonal positivity can fail on non-Delaunay
tetrahedra.  Rotational calibration here means exact affine reconstruction,
not an empirical average of body diagonals.

\section{The shared cochain is the normalized-affine degree}

On a tetrahedron write
\begin{equation}
 v_T(\lambda)=\sum_{i=0}^3\lambda_i n_i,
 \qquad u_T=\frac{v_T}{|v_T|},
 \qquad \lambda\in\Delta^3.
 \label{eq:normalize}
\end{equation}
Fix $-1/3<\epsilon<1$ and impose
$n_i\cdot n_j>\epsilon$ for all vertex pairs of every tetrahedron.  Then
\begin{equation}
 |v_T|^2
 \ge \epsilon+(1-\epsilon)\sum_i\lambda_i^2
 \ge\frac{1+3\epsilon}{4}>0.
 \label{eq:affine-gap}
\end{equation}
The production action uses $\epsilon=-0.32$, for which the last lower bound
is $10^{-2}$; every relaxed solution is in fact much farther from it.

Let $s_T=\operatorname{sgn}\det[x_1-x_0,x_2-x_0,x_3-x_0]$ and define
\begin{equation}
 \Omega_h[n](T)=\frac{s_T}{2\pi^2}\,C_h[n](T)
 \int_{\Delta^3}\frac{\dd\lambda}{|v_T(\lambda)|^4}.
 \label{eq:omega}
\end{equation}

\begin{theorem}[Exact degree identity]
If the boundary is the vacuum and \eqref{eq:affine-gap} holds, the normalized
affine pieces glue to a continuous map $u_h:\Omega/\partial\Omega\to S^3$,
and
\begin{equation}
 \sum_{T\in K_h^{(3)}}\Omega_h[n](T)=\deg u_h\in\ZZ.
 \label{eq:degree}
\end{equation}
\end{theorem}

\begin{proof}
On the reference tetrahedron,
\[
 u_T^*\operatorname{vol}_{S^3}
 =\frac{\det[v_T,\partial_1v_T,\partial_2v_T,\partial_3v_T]}
 {|v_T|^4}\dd\lambda.
\]
The determinant is independent of $\lambda$ and equals
$\det[n_0,n_1,n_2,n_3]=C_h[n](T)$.  Multiplication by $s_T$ converts the
reference orientation to the physical orientation.  Summing and dividing by
$\operatorname{vol}(S^3)=2\pi^2$ is the pullback definition of degree.
\end{proof}

This identity is stronger than agreement with one regular value.  The code
nevertheless checks three regular values and an independent Duffy--Gauss
evaluation of \eqref{eq:omega}.

\section{Compatible complete-minor action and cell formula}

Let $v_h=\Ih n_h$ be the P1 affine interpolant.  The AP-E11 action is
\begin{equation}
 E_h[n_h]=\sum_{T\in K_h}\int_T
 \left\{
 \frac12|Dv_h|^2+\frac R2|\wedge^2Dv_h|^2
 +\frac K2|\wedge^3Dv_h|^2+m^2(1-v_h^0)
 \right\}\dd x,
 \label{eq:action}
\end{equation}
on the admissible nodal $S^3$ sector, and $+\infty$ outside it.  The production
card uses $(R,K,m^2)=(1,0.35,0.12)$.  Equations
\eqref{eq:hodge1}--\eqref{eq:hodge3} express every derivative term as a
positive cochain Hodge norm.  The $K$ term is not an unrelated stabilizer:
it is the norm of the complete family \eqref{eq:three-cup} whose contraction
and radial dressing give \eqref{eq:omega}.  For a smooth $S^3$ map its norm
is the squared topological-current density.

For a periodic P1 corrector $\varphi:Y\to\RR^4$, define
\begin{equation}
 Q_{R,K}^{\rm cell}(A)=\inf_{\varphi\ \mathrm{periodic}}
 \int_Y Q_{R,K}(A+D\varphi)\dd x,
 \quad
 Q_{R,K}(A)=\frac12|A|^2+\frac R2|M_2(A)|^2+\frac K2|M_3(A)|^2.
 \label{eq:cell}
\end{equation}

\begin{theorem}[Finite-$R,K$ mesh-independent cell formula]
On every conforming periodic tetrahedral mesh,
\begin{equation}
 \boxed{Q_{R,K}^{\rm cell}(A)=Q_{R,K}(A),\qquad\varphi=0
 \text{ is a global minimizer}.}
 \label{eq:cell-result}
\end{equation}
This holds for the uniform six-tet mesh and for both checkerboard five-tet
phases, with no mesh-dependent rescaling.
\end{theorem}

\begin{proof}
Periodicity and the null-Lagrangian identities give, for $r=1,2,3$,
\begin{equation}
 \frac1{|Y|}\int_Y M_r(A+D\varphi)\dd x=M_r(A).
 \label{eq:nullmean}
\end{equation}
For $r=1$ this is the mean-zero gradient identity.  For $r=2,3$, expand each
minor and integrate one derivative by parts; periodic boundary terms vanish
and mixed derivatives cancel.  The function
$(F,G,H)\mapsto(|F|^2+R|G|^2+K|H|^2)/2$ is convex.  Jensen applied to the
complete minor vector and \eqref{eq:nullmean} proves the lower bound
$Q_{R,K}(A)$, attained at zero corrector.  This is the standard polyconvex
mechanism \cite{Ball1977}, evaluated here on the actual finite periodic
complex rather than invoked as an asymptotic slogan.
\end{proof}

\section{Compensated compactness and the precise Gamma theorem}

Let $a_h\to0$ on a shape-regular mesh family and suppose
$\sup_hE_h[n_h]<\infty$.  The Dirichlet term bounds $v_h$ in $H^1$.
Rellich gives $v_h\to v$ strongly in $L^2$ after a subsequence.  Nodal unit
length and the local interpolation estimate
\begin{equation}
 \|1-|v_h|\|_{L^2}\le C a_h\|Dv_h\|_{L^2}
 \label{eq:sphere-limit}
\end{equation}
imply $|v|=1$ almost everywhere.

The action bounds $M_2(Dv_h)$ and $M_3(Dv_h)$ in $L^2$.  Because the
cochains are actual affine minors, not one-corner products, the distributional
null-Lagrangian identities identify their weak limits with $M_2(Dv)$ and
$M_3(Dv)$.  Equivalently, each third minor satisfies the Piola form
\begin{equation}
 \det D(v_h^{A},v_h^{B},v_h^{C})
 =\frac13\operatorname{div}
 \left((v_h^{A},v_h^{B},v_h^{C})\,cof
 D(v_h^{A},v_h^{B},v_h^{C})\right).
 \label{eq:piola}
\end{equation}
This is precisely the differential constraint missing in AP-E10; general
div--curl theory explains why such compatibility is indispensable
\cite{BrianeCasadoDiaz2016}.

The normalized current has the exact affine formula
\begin{equation}
 J_h=\frac{\det[v_h,D_1v_h,D_2v_h,D_3v_h]}{|v_h|^4}.
 \label{eq:current}
\end{equation}
The uniform gap \eqref{eq:affine-gap}, strong finite-$L^p$ convergence of
$v_h$, and weak $L^2$ compactness of the third minors identify the weak
limit of $J_h$.  In particular its integral and hence degree are closed.

\begin{theorem}[Unconditional relaxed fixed-degree Gamma theorem]
Fix a degree $B$, the boundary vacuum, $-1/3<\epsilon<1$, and a shape-regular
mesh family.  In strong $L^2$, the discrete functionals \eqref{eq:action}
Gamma-converge on the degree-$B$ sector to the fixed-degree lower-
semicontinuous relaxation of the smooth complete-minor action.
\end{theorem}

\begin{proof}
Equicoercivity and sector closure were established above.  Convex weak lower
semicontinuity in $(Dv,M_2(Dv),M_3(Dv))$ gives the liminf.  For every smooth
$S^3$ map, nodal interpolation is eventually admissible, preserves degree,
and converges in all three Hodge energies by affine exactness.  Apply that
recovery construction to a smooth sequence defining the fixed-degree lower-
semicontinuous relaxation and take a diagonal subsequence.
\end{proof}

The last phrase is important.  This theorem does not assert that the relaxed
functional equals the naive classical graph-norm functional on every finite-
energy Sobolev map.  Smooth approximation of manifold-valued Sobolev maps is
topology-sensitive \cite{Bethuel1991}; the exact density problem for the
simultaneous $Dv$, $M_2(Dv)$, and $M_3(Dv)$ graph norm is not supplied by the
standard theorem.  Recent rigorous lattice-skyrmion work gives a useful
compactness/liminf/recovery architecture in two dimensions
\cite{BrianiCicaleseKreutz2026}, but does not identify this three-dimensional
graph-norm relaxation.

\section{Numerical verification and quotient gate}

The independent card is
\begin{center}
\texttt{route\_f/code/scan\_ap\_e11\_compatible\_cochain\_action.py}.
\end{center}
It verifies exact integer cochain identities, local Hodge matrices, radial
degree integration, periodic cell optimizations, and thirteen unanchored
stationary fields.  Selected residuals are
\begin{center}
\begin{tabular}{@{}lr@{}}
\toprule
check & observed maximum/minimum\\
\midrule
$d^2$, Leibniz, associativity residual & $0$\\
third-cup determinant residual & $8.33\times10^{-16}$\\
minimum local Hodge eigenvalue & $2.485\times10^{-2}$\\
rigid-rotation Hodge residual & $3.55\times10^{-15}$\\
affine one-/two-/three-form reconstruction residual & $2.40\times10^{-16}$\\
radial cochain versus regular-value degree & $3.07\times10^{-7}$\\
periodic complete-minor mean residual & $1.78\times10^{-15}$\\
optimized cell-formula relative residual & $2.56\times10^{-14}$\\
maximum projected gradient density, 13 cases & $1.54\times10^{-7}$\\
minimum pair dot product, 13 cases & $0.2724$\\
\bottomrule
\end{tabular}
\end{center}
Every case has three-target $B=(1,1,1)$.  The common sequence extends through
$N=37$, $L=8$, and $a=0.222222$.  Translation and mesh comparisons use
$N=29$, $L=7$, and $a=0.25$.  The quotient uses the positive density of the same
action, lumped from each exact tetrahedral contribution to its four vertices;
its radial comparison is the maximum CDF distance, not a bin-phase-sensitive
histogram norm.

\begin{center}
\begin{tabular}{@{}lrrr@{}}
\toprule
gate & observed & threshold & pass\\
\midrule
joint-tail energy spread & $0.437\%$ & $3\%$ & yes\\
joint-tail centered-radius spread & $3.489\%$ & $5\%$ & yes\\
joint-tail radial CDF distance & $4.067\%$ & $5\%$ & yes\\
translated-start energy spread & $8.33\times10^{-12}$ & $1\%$ & yes\\
translated-start radius spread & $7.16\times10^{-7}$ & $2\%$ & yes\\
translated-start CDF distance & $5.91\times10^{-7}$ & $3\%$ & yes\\
five-/six-tet energy spread & $0.463\%$ & $3\%$ & yes\\
five-/six-tet radius spread & $0.752\%$ & $5\%$ & yes\\
five-/six-tet CDF distance & $3.556\%$ & $5\%$ & yes\\
\bottomrule
\end{tabular}
\end{center}

Thus
\begin{equation}
 \boxed{C_{\rm cell}=\mathrm{true},\quad
 C_{\rm quotient}=\mathrm{true},\quad
 C_{\rm relaxed\ regulator}=\mathrm{true}.}
 \label{eq:passed}
\end{equation}
The classical-action identification remains
\begin{equation}
 \boxed{C_{\rm graph\ density}=\mathrm{false},\quad
 C_{\rm classical\ regulator}=\mathrm{false},\quad
 C_{\rm Hessian}=C_{\det}=\mathrm{false}.}
 \label{eq:failed}
\end{equation}
Equation \eqref{eq:failed} is not a return of the AP-E10 stencil blocker.
That blocker is removed.  It is a sharper functional-analytic question:
either prove fixed-degree smooth density in the complete-minor graph norm, or
prove that the selected relaxed minimizer is a regular isolated classical
solution.  Only then may the same-action Riemann Hessian be assembled; a
regulated determinant additionally requires a declared fermion kernel and
counterterm prescription.

\bibliographystyle{unsrt}
\bibliography{route_f/tex/ap_e11_compatible_cochain_action}

\end{document}
